\documentclass[11pt]{amsart}
\usepackage[margin=1in]{geometry}
\usepackage{amsmath,amssymb,amsthm,mathtools}
\usepackage{hyperref}
\usepackage{graphicx}

\newtheorem{theorem}{Theorem}
\newtheorem{lemma}{Lemma}
\newtheorem{corollary}{Corollary}
\theoremstyle{remark}
\newtheorem{remark}{Remark}

\newcommand{\G}{\mathbb G}
\newcommand{\hG}{\widehat{\mathbb G}}
\newcommand{\TC}{\mathcal T}
\newcommand{\B}{\mathcal B}
\newcommand{\K}{\mathcal K}
\newcommand{\cb}{\mathrm{cb}}
\newcommand{\id}{\mathrm{id}}
\newcommand{\opmod}{\mathbf{mod}}
\newcommand{\hot}{\widehat{\otimes}}
\newcommand{\vnot}{\overline{\otimes}}
\newcommand{\rhdp}{\triangleright}
\newcommand{\lhdp}{\triangleleft}

\title[Trace-Class Projectivity Replacement]{Replacing Relative Projectivity by Projectivity for the Trace-Class Convolution Module}
\author{Agent lane 17}
\date{2026-07-18}

\begin{document}

\begin{abstract}
We answer Remark 5.14 of Crann's paper \emph{Amenability and covariant injectivity of locally compact quantum groups II}.  For the trace-class convolution modules in Theorems 5.12 and 5.13, relative 1-projectivity may be replaced by ordinary metric 1-projectivity.  More precisely, for every locally compact quantum group \(\G\), compactness of \(\hG\), relative 1-projectivity of \((\TC(L^2(\G)),\rhdp)\) as a left module over itself, and 1-projectivity of the same module are equivalent.  The analogous right-module statement for \((\TC(L^2(\G)),\lhdp)\) also holds.
\end{abstract}

\maketitle

\section{Source Question}

The source paper proves the following theorem, in its notation \cite[Theorems 5.12 and 5.13]{Crann2015}.  For a locally compact quantum group \(\G\), the conditions
\[
 \hG\text{ compact},\qquad
 \TC(L^2(\G))\text{ relatively 1-projective in }(\TC(L^2(\G)),\rhdp)-\opmod
\]
are equivalent; there is also a right-module analogue for \((\TC(L^2(\G)),\lhdp)\).  Immediately after those theorems, Remark 5.14 asks whether the word ``relative'' can be removed from condition (2) \cite[Remark 5.14]{Crann2015}.

\begin{center}
\includegraphics[width=.92\linewidth]{figures/open_problem_crop.png}
\end{center}

\section{Proof Intuition}

There are two different reasons why the replacement works.

First, ordinary 1-projectivity is strong enough to recover the relative splitting.  Apply projectivity to the multiplication
\[
 m=(\Gamma^r)_*: \TC(L^2(\G))\hot \TC(L^2(\G))\longrightarrow \TC(L^2(\G)).
\]
This gives exact splittings whose norms are \(<1+\varepsilon\).  Dualizing puts us inside the dual operator space of normal completely bounded maps
\[
 \B(L^2(\G)\otimes L^2(\G))\longrightarrow \B(L^2(\G)).
\]
A weak-star cluster point as \(\varepsilon\to0\) has norm at most 1 and is still a module left inverse of \(\Gamma^r\).  Passing back to the predual gives the norm-one relative projective splitting.

Second, the reverse implication uses the special trace-class convolution algebra rather than plain operator-space projectivity.  If \(\hG\) is compact, then \(\G\) is discrete.  Hu, Neufang, and Ruan show that in this case \((\TC(L^2(\G)),\rhdp)\) has a right identity obtained by extending the identity of \(L^1(\G)\).  A completely contractive Banach algebra with a contractive right identity is automatically 1-projective as a left module over itself.

\section{The Result}

\begin{theorem}\label{t:main-left}
Let \(\G\) be a locally compact quantum group and set
\[
 A_r=(\TC(L^2(\G)),\rhdp).
\]
The following are equivalent.
\begin{enumerate}
\item \(\hG\) is compact, equivalently \(\G\) is discrete.
\item \(A_r\) is relatively 1-projective in \(A_r-\opmod\).
\item \(A_r\) is 1-projective in \(A_r-\opmod\).
\end{enumerate}
\end{theorem}

\begin{theorem}\label{t:main-right}
Let
\[
 A_l=(\TC(L^2(\G)),\lhdp).
\]
The following are equivalent.
\begin{enumerate}
\item \(\hG\) is compact, equivalently \(\G\) is discrete.
\item \(A_l\) is relatively 1-projective in \(\opmod-A_l\).
\item \(A_l\) is 1-projective in \(\opmod-A_l\).
\end{enumerate}
\end{theorem}

\begin{lemma}\label{l:identity-projective}
Let \(A\) be a completely contractive Banach algebra.
\begin{enumerate}
\item If \(A\) has a right identity \(e\) with \(\|e\|\leq1\), then \(A\) is 1-projective in \(A-\opmod\).
\item If \(A\) has a left identity \(e\) with \(\|e\|\leq1\), then \(A\) is 1-projective in \(\opmod-A\).
\end{enumerate}
\end{lemma}

\begin{proof}
We prove the left-module statement; the right-module statement is identical with sides reversed.  Let \(q:Y\twoheadrightarrow Z\) be a complete quotient morphism in \(A-\opmod\), and let \(\phi:A\to Z\) be a completely bounded \(A\)-module morphism.  Given \(\varepsilon>0\), choose \(y\in Y\) such that
\[
 q(y)=\phi(e),\qquad \|y\|<\|\phi(e)\|+\varepsilon\leq \|\phi\|_{\cb}+\varepsilon.
\]
Define \(\widetilde\phi(a)=a\cdot y\).  The module action is completely contractive, so \(\|\widetilde\phi\|_{\cb}\leq\|y\|\).  Moreover
\[
 q\widetilde\phi(a)=a\cdot q(y)=a\cdot\phi(e)=\phi(ae)=\phi(a),
\]
since \(e\) is a right identity.  Thus \(A\) satisfies the metric lifting property with constant 1.
\end{proof}

\begin{lemma}\label{l:projective-relative}
Let \(A_r=(\TC(L^2(\G)),\rhdp)\).  If \(A_r\) is 1-projective in \(A_r-\opmod\), then \(A_r\) is relatively 1-projective in \(A_r-\opmod\).
\end{lemma}

\begin{proof}
Let \(H=L^2(\G)\).  The right fundamental unitary gives a normal complete isometry
\[
 \Gamma^r:\B(H)\to \B(H)\vnot\B(H),
\]
and the multiplication \(m=(\Gamma^r)_*:A_r\hot A_r\to A_r\) is a complete quotient morphism in \(A_r-\opmod\).  Since \(A_r\) is 1-projective, for each \(n\geq1\) there is an \(A_r\)-module morphism
\[
 \sigma_n:A_r\to A_r\hot A_r
\]
such that \(m\sigma_n=\id_{A_r}\) and \(\|\sigma_n\|_{\cb}<1+1/n\).

Dualizing, \(S_n=\sigma_n^*\) is a normal completely bounded module map
\[
 S_n:\B(H)\vnot\B(H)\to \B(H)
\]
with \(S_n\Gamma^r=\id_{\B(H)}\) and \(\|S_n\|_{\cb}<1+1/n\).  The space of normal completely bounded maps \(\mathcal{CB}^{\sigma}(\B(H)\vnot\B(H),\B(H))\) is a dual operator space.  Hence a subnet of \((S_n)\) has a weak-star cluster point \(S\) of cb-norm at most 1.

The equations defining \(S_n\Gamma^r=\id\) and the module-intertwining identities are weak-star closed: after pairing with elements of the predual \(A_r\), they are linear equations.  Therefore \(S\Gamma^r=\id_{\B(H)}\), and \(S\) is still the corresponding normal module morphism.  Since \(S\) is normal, it has a preadjoint
\[
 \sigma=S_*:A_r\to A_r\hot A_r.
\]
The relations above give \(m\sigma=\id_{A_r}\), \(\|\sigma\|_{\cb}\leq1\), and \(\sigma\) is an \(A_r\)-module morphism.  The density identity
\[
 A_r=\overline{\operatorname{span}}(A_r\rhdp A_r)
\]
is proved in the source paper, so this is exactly relative 1-projectivity in the non-unitized form used there.
\end{proof}

\begin{proof}[Proof of Theorem \ref{t:main-left}]
The equivalence of (1) and (2) is Theorem 5.12 of the source paper \cite[Theorem 5.12]{Crann2015}.

If (3) holds, Lemma \ref{l:projective-relative} gives (2), and hence (1).

Conversely assume (1).  Then \(\G\) is discrete.  Hu, Neufang, and Ruan show that if \(f_0\) is the identity of \(L^1(\G)\) and \(\omega_0\in\TC(L^2(\G))\) is a normal extension of \(f_0\), then \(\omega_0\) is a right identity for \((\TC(L^2(\G)),\rhdp)\) \cite[Proposition 5.1(iii)]{HNR2013}.  The discrete counit \(f_0\) is a normal state, and the normal Hahn--Banach extension theorem gives a norm-one normal extension \(\omega_0\) to \(\B(L^2(\G))\).  Lemma \ref{l:identity-projective} therefore gives 1-projectivity of \(A_r\) in \(A_r-\opmod\).  Thus (1) implies (3).
\end{proof}

\begin{proof}[Proof of Theorem \ref{t:main-right}]
The source paper's Theorem 5.13 gives (1) \(\Leftrightarrow\) relative 1-projectivity of \(A_l\) in \(\opmod-A_l\) \cite[Theorem 5.13]{Crann2015}.  The proof of Lemma \ref{l:projective-relative} applies verbatim to the left fundamental unitary and the multiplication \(\lhdp\), with right modules in place of left modules, giving ordinary 1-projectivity \(\Rightarrow\) relative 1-projectivity.

For the reverse implication, assume \(\G\) is discrete.  The unitary antipode induces a complete isometric anti-isomorphism \(R_*:A_r\to A_l\) satisfying
\[
 R_*(\omega\rhdp\tau)=R_*(\tau)\lhdp R_*(\omega).
\]
If \(\omega_0\) is the contractive right identity of \(A_r\), then \(R_*(\omega_0)\) is a contractive left identity of \(A_l\).  Lemma \ref{l:identity-projective} gives 1-projectivity in \(\opmod-A_l\).
\end{proof}

\section{Verification and Novelty Notes}

The argument is formal after two external inputs:
\begin{enumerate}
\item Crann's Theorems 5.12 and 5.13 \cite{Crann2015}, which identify relative 1-projectivity of the two trace-class modules with compactness of \(\hG\).
\item Hu-Neufang-Ruan's Proposition 5.1(iii) \cite{HNR2013}, which says that when \(\G\) is discrete, a normal extension of the identity of \(L^1(\G)\) is a right identity for the trace-class convolution algebra \((\TC(L^2(\G)),\rhdp)\).
\end{enumerate}

The local run indexes were searched for the arXiv id, the title, ``relative 1-projectivity'', ``TCQ'', and trace-class projectivity; no duplicate settlement was found.  Web and local source searches found the source paper, the HNR trace-class convolution paper, Crann's 2021 local-theory paper, and later trace-class convolution papers, but no explicit answer to Remark 5.14.

\section{References}

\begin{thebibliography}{9}

\bibitem{Crann2015}
J. Crann,
\emph{Amenability and covariant injectivity of locally compact quantum groups II},
Canadian Journal of Mathematics 69 (2017), 1064--1086.
arXiv:1507.03296.

\bibitem{HNR2013}
Z. Hu, M. Neufang, and Z.-J. Ruan,
\emph{Convolution of trace class operators over locally compact quantum groups},
Canadian Journal of Mathematics 65 (2013), no. 5, 1043--1072.
doi:10.4153/CJM-2012-030-5.

\end{thebibliography}

\end{document}
